% Source: https://artofproblemsolving.com/community/c3h1132278_odd_perfect_numbers
% Stable identity: aops-topic-1132278-q01
The first sentence is not a known theorem. The Euclid--Euler theorem says that every \emph{even} perfect number has the form $2^{p-1}(2^p-1)$ with $2^p-1$ prime (Euclid showed that such numbers are perfect, and Euler showed that every even perfect number is of this form). Euler's half of the proof uses the fact that the number is even: it writes $n=2^{k}m$ with $m$ odd and $k\ge1$.

So the statement ``every perfect number has this form'' already assumes that the number is even, and the argument is circular: it proves ``every even perfect number is even''. Whether an odd perfect number exists is a famous unsolved problem; the argument says nothing about it.
